% Examen Parcial CS4054 2025-I, recompuesto en LaTeX a partir de una copia
% impresa (sin datos de alumno ni respuestas). Compilar:
%   cd content/evaluaciones/parciales/2025-i && latexmk -pdf ep-2025-i.tex
%   cp ep-2025-i.pdf ../../../../public/evaluaciones/parciales/2025-i.pdf
\documentclass[10pt,a4paper]{article}
\usepackage[T1]{fontenc}
\usepackage[utf8]{inputenc}
\usepackage[spanish]{babel}
\usepackage{lmodern}
\usepackage[margin=1.8cm,top=1.6cm]{geometry}
\usepackage{amsmath,amssymb}
\usepackage{tabularx,booktabs,array}
\usepackage{enumitem}
\usepackage{fancyvrb}
\usepackage{xcolor}
\usepackage{tikz}
\usetikzlibrary{arrows.meta}
\usepackage{fancyhdr}
\usepackage[hidelinks]{hyperref}

\definecolor{wsdark}{HTML}{2B2B2B}
\definecolor{wsrule}{HTML}{B0B0B0}
\setlength{\parindent}{0pt}
\setlength{\parskip}{4pt}
\setlist{itemsep=2pt,topsep=3pt}

\pagestyle{fancy}\fancyhf{}
\renewcommand{\headrulewidth}{0pt}
\fancyhead[L]{\small\{garias,cwilliams\}@utec.edu.pe\\\textbf{CS4055: Networks and communications}\\2025.05.23}
\fancyhead[R]{\textbf{\large UTEC}}
\setlength{\headheight}{36pt}
\fancyfoot[C]{\small\thepage}

\newcommand{\pts}[1]{\textbf{(#1)}}
\newenvironment{trace}{\par\small\VerbatimEnvironment\begin{Verbatim}[fontsize=\scriptsize,samepage=true,frame=single,rulecolor=\color{wsrule},framesep=2mm]}{\end{Verbatim}}

\begin{document}

\begin{center}
  {\Large\bfseries CS4054 -- EXAMEN PARCIAL 2025-I}\\[6pt]
  \renewcommand{\arraystretch}{1.6}
  \begin{tabularx}{\linewidth}{|>{\bfseries}l|X|>{\bfseries}l|p{1.6cm}|}
    \hline
    Apellidos y Nombres: & & Nota: & \\ \hline
  \end{tabularx}
\end{center}

\textbf{Instrucciones:}
\begin{itemize}\small\itshape
  \item Realizar la evaluación de manera \textbf{individual}, utilizando \textbf{lapicero de color azul o negro (no se revisará escritura con lápiz).}
  \item Tomar en cuenta que se calificará el procedimiento \textbf{(respuesta final sin procedimiento no será tomado en cuenta).}
  \item Sea cuidadoso con su \textbf{redacción}, la cual formará parte de su calificación (aspectos léxicos, sintácticos, semánticos).
  \item Lea bien la pregunta o enunciado antes de proceder a su desarrollo, \textbf{administre su tiempo eficazmente.}
  \item \textbf{Sólo se revisará el procedimiento evidenciado en su cuadernillo de resolución.} No se tomará en cuenta procedimiento escrito sobre la hoja de evaluación o sobre hojas borradoras.
  \item Pregunta cuya solución \textbf{no sea legible, NO SERÁ REVISADA.}
\end{itemize}

\section*{A. Sección Teórica (11 pts.)}
\begin{enumerate}[label=\arabic*.]
\item \pts{1.5 pt} Explique usted las diferencias entre TCP y UDP (mencione y explique tres diferencias). Luego, supongamos que quiere establecer comunicación desde un cliente remoto a un servidor lo más rápido posible. ¿Lo más adecuado sería utilizar UDP o TCP? ¿Por qué?

\item \pts{3.0 pt} Considere un archivo de tamaño $F$ y $N$ peers. Un servidor con velocidad de subida $u_s$ intenta transmitir el archivo hacia los $N$ peers, cada uno con su velocidad de descarga $d_n$ y velocidad de subida $u_n$. Además $d_{min} = \min\{d_1, d_2, \dots, d_N\}$ es la velocidad de descarga mínima entre los peers. Explique cómo se obtienen las siguientes expresiones del tiempo de distribución total para las arquitecturas cliente-servidor y P2P e indique las implicancias cuando el número de peers es muy alto
\[
  D_{cs} \ge \max\left\{\frac{NF}{u_s}, \frac{F}{d_{\min}}\right\}.
  \qquad\qquad
  D_{P2P} \ge \max\left\{\frac{F}{u_s}, \frac{F}{d_{\min}}, \frac{NF}{u_s + \sum_{i=1}^{N} u_i}\right\}
\]

\item \pts{2.5 pt} (Capa de Transporte) Identifique y justifique el uso de los protocolos de transporte utilizados por algunas aplicaciones populares de Internet. \textbf{(Llenar EN SU CUADERNILLO)}

\smallskip
\renewcommand{\arraystretch}{1.25}
\begin{tabularx}{\linewidth}{|l|c|c|X|}
  \hline
  \textbf{Aplicación} & \textbf{Protocolo de aplicación} & \textbf{Protocolo de transporte} & \textbf{Justificación de protocolo de transporte} \\ \hline
  Email & SMTP & & \\ \hline
  Acceso terminal remoto & Telnet & & \\ \hline
  Web & HTTP 1.1 & & \\ \hline
  Transferencia de archivos & FTP & & \\ \hline
  Transmisión multimedia & HTTP, DASH & & \\ \hline
  DNS & DNS & & \\ \hline
\end{tabularx}

\item \pts{2.0 pt} Resolver el \textbf{FSM de un RDT 2.1} (reliable data transfer). Acompañe su respuesta con una explicación de cómo funciona.

\item \pts{2.0 pt} Explicar el cambio que haría para poder pasar de \textbf{RDT 2.1} a \textbf{RDT 2.2}.
\end{enumerate}

\section*{B. Sección Práctica (9 pts.)}
\begin{enumerate}[label=\arabic*.,start=6]
\item \pts{1 pt.} Se muestra el siguiente resultado de ejecución del comando nslookup con el filtrado del tipo: \textbf{``set type=mx''}. ¿Qué servidor se contactará primero cuando se envíe correo electrónico a la dirección propuesta? ¿Por qué?
\begin{trace}
> netflix.com
Server:  UnKnown
Address:  192.168.107.95

Non-authoritative answer:
netflix.com     MX preference = 1, mail exchanger = aspmx.l.google.com
netflix.com     MX preference = 10, mail exchanger = aspmx2.googlemail.com
netflix.com     MX preference = 10, mail exchanger = aspmx3.googlemail.com
netflix.com     MX preference = 5, mail exchanger = alt1.aspmx.l.google.com
netflix.com     MX preference = 5, mail exchanger = alt2.aspmx.l.google.com
\end{trace}

\item \pts{3.0 pts} \textbf{Llene la siguiente tabla en su cuadernillo para 5 segmentos}, referente a los tiempos RTT (vista en el laboratorio), empleando los datos de los traces de las \textbf{Figura 1 y figura 2.}

Emplee la fórmula: \quad $\mathit{EstimatedRTT} = 0.875 \cdot \mathit{EstimatedRTT} + 0.125 \cdot \mathit{SampleRTT}$
\begin{trace}
No. Time      Source          Destination     Protocol Info
 1  0.000000  192.168.1.102   128.119.245.12  TCP   1161 > http [SYN] Seq=0 Ack=0 Win=16384 Len=0 MSS=1460
 2  0.023172  128.119.245.12  192.168.1.102   TCP   http > 1161 [SYN, ACK] Seq=0 Ack=1 Win=5840 Len=0 MSS=1460
 3  0.023265  192.168.1.102   128.119.245.12  TCP   1161 > http [ACK] Seq=1 Ack=1 Win=17520 Len=0
 4  0.026477  192.168.1.102   128.119.245.12  HTTP  POST /ethereal-labs/lab3-1-reply.htm HTTP/1.1
 5  0.041737  192.168.1.102   128.119.245.12  HTTP  Continuation or non-HTTP traffic
 6  0.053937  128.119.245.12  192.168.1.102   TCP   http > 1161 [ACK] Seq=1 Ack=566 Win=6780 Len=0
 7  0.054026  192.168.1.102   128.119.245.12  HTTP  Continuation or non-HTTP traffic
 8  0.054690  192.168.1.102   128.119.245.12  HTTP  Continuation or non-HTTP traffic
 9  0.077294  128.119.245.12  192.168.1.102   TCP   http > 1161 [ACK] Seq=1 Ack=2026 Win=8760 Len=0
10  0.077405  192.168.1.102   128.119.245.12  HTTP  Continuation or non-HTTP traffic
11  0.078157  192.168.1.102   128.119.245.12  HTTP  Continuation or non-HTTP traffic
12  0.124085  128.119.245.12  192.168.1.102   TCP   http > 1161 [ACK] Seq=1 Ack=3486 Win=11680 Len=0
13  0.124185  192.168.1.102   128.119.245.12  HTTP  Continuation or non-HTTP traffic
\end{trace}
\vspace{-6pt}\begin{center}\small\textbf{Figura 1.} POST y Continuations.\end{center}
\begin{trace}
No. Time      Source          Destination     Protocol Info
 5  0.041737  192.168.1.102   128.119.245.12  HTTP  Continuation or non-HTTP traffic
 6  0.053937  128.119.245.12  192.168.1.102   TCP   http > 1161 [ACK] Seq=1 Ack=566 Win=6780 Len=0
 7  0.054026  192.168.1.102   128.119.245.12  HTTP  Continuation or non-HTTP traffic
 8  0.054690  192.168.1.102   128.119.245.12  HTTP  Continuation or non-HTTP traffic
 9  0.077294  128.119.245.12  192.168.1.102   TCP   http > 1161 [ACK] Seq=1 Ack=2026 Win=8760 Len=0
10  0.077405  192.168.1.102   128.119.245.12  HTTP  Continuation or non-HTTP traffic
11  0.078157  192.168.1.102   128.119.245.12  HTTP  Continuation or non-HTTP traffic
12  0.124085  128.119.245.12  192.168.1.102   TCP   http > 1161 [ACK] Seq=1 Ack=3486 Win=11680 Len=0
13  0.124185  192.168.1.102   128.119.245.12  HTTP  Continuation or non-HTTP traffic
14  0.169118  128.119.245.12  192.168.1.102   TCP   http > 1161 [ACK] Seq=1 Ack=4946 Win=14600 Len=0
15  0.217299  128.119.245.12  192.168.1.102   TCP   http > 1161 [ACK] Seq=1 Ack=6406 Win=17520 Len=0
16  ...                                       TCP   http > 1161 [ACK] ...
17  0.304807  128.119.245.12  192.168.1.102   TCP   http > 1161 [ACK] Seq=1 Ack=9013 Win=23360 Len=0
\end{trace}
\vspace{-6pt}\begin{center}\small\textbf{Figura 2.} Acknowledgments correspondientes.\end{center}

\renewcommand{\arraystretch}{1.3}
\begin{tabularx}{\linewidth}{|X|X|X|X|X|}
  \hline
  \centering N° & \centering Sent time & \centering ACK Received time & \centering RTT (seconds) & \centering\arraybackslash EstimatedRTT \\ \hline
  \centering Seg1 \dots & & & & \\ \hline
\end{tabularx}

\item Basado en el siguiente diagrama de dispositivos, con sus diferentes direcciones IP asignadas. Suponga que quiere realizar una comunicación UDP utilizando Python.

\begin{center}
\begin{tikzpicture}[font=\small, x=1cm, y=1cm]
  \newcommand{\lap}[1]{%
    \begin{scope}[shift={#1}]
      \draw[fill=gray!30, rounded corners=0.6pt] (-0.42,0.06) rectangle (0.42,0.58);
      \fill[gray!65] (-0.36,0.11) rectangle (0.36,0.53);
      \draw[fill=gray!45] (-0.5,0.06) -- (0.5,0.06) -- (0.6,-0.06) -- (-0.6,-0.06) -- cycle;
    \end{scope}}
  \draw[wsrule] (-0.6,-1.1) rectangle (10.4,3.3);
  \lap{(4.9,1.9)} \lap{(1.6,0)} \lap{(7.9,0)}
  \node[font=\small\bfseries] at (4.9,2.85) {192.168.1.3};
  \node[align=center] at (4.9,1.35) {Laptop-PT\\Alumno B};
  \node[font=\small\bfseries] at (1.6,0.95) {192.168.1.4};
  \node[align=center] at (1.6,-0.55) {Laptop-PT\\Alumno A};
  \node[font=\small\bfseries] at (9.25,0.75) {192.168.1.10};
  \node[align=center] at (7.9,-0.55) {Laptop-PT\\Alumno C};
  \draw[-{Stealth[length=2.6mm]}, thick, red!70!black] (2.4,0.25) -- (7.1,0.25);
  \draw[-{Stealth[length=2.6mm]}, thick, red!70!black] (5.6,2.05) -- (7.5,0.72);
\end{tikzpicture}
\end{center}
\begin{enumerate}[label=\alph*.]
  \item \pts{1 pt.} Si se quiere enviar un mensaje del alumno A a la computadora C, qué datos necesita indicar en el datagrama UDP de su función \texttt{sender.py}.
  \item \pts{1 pt.} Si se quiere recibir al mismo tiempo, los mensajes del alumno B en la laptop C, ¿cuál es la configuración que se debe realizar? ¿Qué consideraciones se deben tener en cuenta?
  \item \pts{0.5 pt.} Cuál es el tamaño máximo de puerto en UDP.
\end{enumerate}

\item Basado en los siguientes datos obtenidos de un frame de comunicación en Wireshark. Responda:
\begin{trace}
v Internet Protocol Version 4, Src: 192.168.107.95, Dst: 192.168.107.165
    0100 .... = Version: 4
    .... 0101 = Header Length: 20 bytes (5)
  > Differentiated Services Field: 0x00 (DSCP: CS0, ECN: Not-ECT)
    Total Length: 152
    Identification: 0x97fb (38907)
  > 010. .... = Flags: 0x2, Don't fragment
    ...0 0000 0000 0000 = Fragment Offset: 0
    Time to Live: 64
    Protocol: UDP (17)
    Header Checksum: 0x4a04 [validation disabled]
    [Header checksum status: Unverified]
    Source Address: 192.168.107.95
    Destination Address: 192.168.107.165
    [Stream index: 8]
v User Datagram Protocol, Src Port: 53, Dst Port: 59270
    Source Port: 53
    Destination Port: 59270
    Length: 132
    Checksum: 0x41a2 [unverified]
    [Checksum Status: Unverified]
    [Stream index: 1]
    [Stream Packet Number: 2]
  > [Timestamps]
    UDP payload (124 bytes)
v Domain Name System (response)
    Transaction ID: 0x000c
  > Flags: 0x8180 Standard query response, No error
    Questions: 1
    Answer RRs: 6
    Authority RRs: 0
    Additional RRs: 0
  > Queries
  > Answers
    [Time: 0.021096000 seconds]
\end{trace}
\begin{enumerate}[label=\alph*.]
  \item \pts{0.5 pt.} ¿Cuáles son las direcciones IP de origen y destino?
  \item \pts{0.5 pt.} ¿Cuál es la longitud del payload (en bytes)?
  \item \pts{0.5 pt.} ¿Qué implicaciones tiene que el checksum aparezca como ``unverified''?
  \item \pts{1 pt.} Brinde el datagrama UDP correspondiente.
\end{enumerate}
\end{enumerate}

\end{document}
